% Generated draft: requires mathematical and source review.
\documentclass[11pt]{article}
\usepackage[margin=25mm]{geometry}
\usepackage{fontspec}
\setmainfont{Latin Modern Roman}
\usepackage{amsmath,amssymb,mathtools,unicode-math}
\setmathfont{Latin Modern Math}
\usepackage{enumitem,needspace,xurl,hyperref}
\setlength{\parindent}{0pt}
\setlength{\parskip}{5pt}
\setlength{\emergencystretch}{3em}
\newcommand{\N}{\mathbb{N}}
\newcommand{\Z}{\mathbb{Z}}
\newcommand{\Q}{\mathbb{Q}}
\newcommand{\R}{\mathbb{R}}
\newcommand{\C}{\mathbb{C}}
\newcommand{\F}{\mathbb{F}}
\newcommand{\A}{\mathbb{A}}
\newcommand{\PP}{\mathbb{P}}
\newcommand{\E}{\mathbb{E}}
\newcommand{\eps}{\varepsilon}
\newcommand{\dd}{\,\mathrm{d}}
\begin{document}
\section*{UnsolvedMath ID: 11000240}
\textbf{Source problem code:} catalogue entry
\subsection{Definitions and mathematical statement}
For the source's \(G\leq\operatorname{MCG}(S)\), retain its surface-group extension
\[
 1\longrightarrow\pi_1(S)\longrightarrow\Gamma_G\longrightarrow G\longrightarrow1.
\]
The quasi-isometry group \(\operatorname{QI}(\Gamma_G)\) consists of self quasi-isometries of a word metric modulo finite uniform distance. Let \(\operatorname{QI}_f(\Gamma_G)\) denote the subgroup represented by maps coarsely preserving the family of cosets of \(\pi_1(S)\): the coset images, and those under a coarse inverse, lie at uniformly bounded Hausdorff distance from cosets. The exact source convention is not included.

The task is to explore or compute the source's group \(\Gamma_C\) for other free subgroups \(G\), including subgroups generated by high powers of twists about a pair of filling curves. Filling means every essential curve meets the pair, with the surface's boundary/puncture convention fixed. The construction of \(C\) and \(\Gamma_C\) is absent, so \(\Gamma_C\) is not identified with an invented commensurator or extension.

The continuation explains that [Mos03a] gives a generalization of the missing Theorem 5.3 for nonfree \(G\), involving \(\operatorname{QI}_f\) rather than all quasi-isometries. It reports coarse algebraic-topology methods [FM00] giving \(\operatorname{QI}_f=\operatorname{QI}\) in the free case, and Whyte's homological methods giving this equality for \(G=\operatorname{MCG}(S)\) in Theorem 5.2. It suggests extending the latter methods to Bieri--Eckmann duality groups that are not Poincar\'e duality groups, with the mapping class group cited as an example [Har86]. The exact theorem hypotheses are missing.

For a Leininger--Reid subgroup, reported here to be a two-dimensional Poincar\'e duality group, the source says those methods fail and nevertheless asks to calculate
\[
 \Gamma_C=\operatorname{QI}_f(\Gamma_G).
\]
This final target does not assert that every quasi-isometry preserves fibers or that no other method can overcome the reported obstacle.
\subsection{Short English statement}
Compute the indicated quasi-isometric/extension group for further free subgroups and for Leininger--Reid surface subgroups, and investigate when all quasi-isometries preserve the surface-group cosets.
\subsection{Sources}
\begin{itemize}
\item ``Problems in the geometry of surface group extensions,'' in \emph{Problems on Mapping Class Groups and Related Topics}, edited by Farb, Problem 5.5 and continuation, supplied PDF page 275. \url{https://www.math.uchicago.edu/~farb/papers/mcgbook.pdf}. Recorded access 2026-07-29; extraction pdf-font-structure.
\item Source [Mos03a], [FM00] and [Har86]; Whyte's homological methods; Leininger--Reid subgroup construction. Full bibliography, Theorems 5.2/5.3 and the construction of \(\Gamma_C\) are absent.
\item Local UnsolvedMath registry, \texttt{data/unsolvedmath/problems.json}, permanent ID \(11000240\), exact source code \texttt{catalogue entry}; supplied OPEN-TRIAGE summary reports unverified status.
\end{itemize}
\end{document}
